\section{Part~2: residual data elements and equivariant nominal models}
\label{sec:part2}

Part~0 stated $H_0$ for the raw data element $Z_k$; Part~1 (core,
Section~\ref{sec:part1core}) gives the exact test and the isotypic reading.
Sprint~4 introduced a second reading of the same telemetry: a nominal model
$\hat f$ (the analytic generalized-momentum observer
\cite{deluca2003actuator,haddadin2017robot,evangelisti2024datadriven}, or a
learned Deep Lagrangian Network \cite{lutter2019delan}) and the residual
$r = y - \hat f(x,u)$, read by a magnitude channel $R^{+}$ (Block~L1). Until
now the two channels were asymmetric: $R^{-}$ tested the raw signal, $R^{+}$
the residual. This part removes the asymmetry. Its content is: (\S\ref{sec:p2-defs})
what an \emph{equivariant} nominal model is and how its defect is measured;
(\S\ref{sec:p2-n31}) the residual inherits $H_0$ from the data, so
Theorems~\ref{thm:n1-1} and~\ref{thm:n1-2} apply to the residual data element
verbatim, and a non-equivariant model contaminates $H_0$ by a fifth term
$\epsbar{model}$ in the $\varepsilon$-budget; (\S\ref{sec:p2-n32}) when the
residual channel has more power than the raw channel; (\S\ref{sec:p2-learning})
where the learned model sits --- the calibration-centring lemma, and an
equivariant DeLaN obtained by mirror weight sharing, with $\delta_f = 0$ and
the Lagrangian structure intact; (\S\ref{sec:p2-budget}) the updated
$\varepsilon$-budget table and the three-channel architecture. Every
proposition carries the Part~0 triple: proof (or complete skeleton),
falsification condition, empirical anchor. Anchors point to the Sprint~6
experiment family e13 (Section~\ref{sec:p2-anchors}); the numbers are
back-filled from the e13 run.

\subsection{Equivariant nominal models and the residual data element}
\label{sec:p2-defs}

\paragraph{Nominal model.}
A nominal model is a map $\hat f : X \times U \to Y_r$ from (state, input)
to a torque-valued output space $Y_r$ (for the Go2/M1: $Y_r = \R^{n_j}$, one
row per actuated joint; the momentum observer adds six floating-base rows,
Section~\ref{sec:p2-budget}). The measured counterpart of $\hat f$ is a
known signal $y \in Y_r$ built from the telemetry --- for the joint rows,
$y = \tau_{\mathrm{cmd}} + J(q)^{\top} f_{\mathrm{c}}$, commanded torque
plus the contact-wrench term (Sprint~4, Finding~1: the \emph{commanded}
torque is the right source; the measured output torque makes the residual
physically consistent and blind to actuator gain/bias faults). The
representation of $G$ on $Y_r$ is the torque representation,
$\repR := \repU$ on the joint block (torques are conjugate to joint
coordinates, Section~\ref{sec:reps}), extended by the polar/axial rules on
any base rows.

\begin{definition}[$G$-equivariant nominal model]
\label{def:equiv-model}
$\hat f$ is $G$-equivariant if
\begin{equation}
  \label{eq:equiv-model}
  \hat f\bigl(\repX(g)\,x,\; \repU(g)\,u\bigr) \;=\; \repR(g)\, \hat f(x,u)
  \qquad \forall g \in G,\ (x,u) \in X \times U .
\end{equation}
\end{definition}

\begin{definition}[Equivariance defect]
\label{def:defect}
The (sup) defect of $\hat f$ on the operating region $\mathcal{K}$ is
\begin{equation}
  \label{eq:defect}
  \delta_f \;:=\; \sup_{g \in G,\ (x,u) \in \mathcal{K}}
  \bigl\| \hat f(\repX(g)x, \repU(g)u) - \repR(g)\hat f(x,u) \bigr\| ,
\end{equation}
and, writing $\delta_f(x,u) := \max_g \|\hat f(\repX(g)x,\repU(g)u) -
\repR(g)\hat f(x,u)\|$ for the pointwise defect, the \emph{quantile defect}
under the nominal per-sample law $P_{\mathrm{nom}}$ of $(x,u)$ is
\begin{equation}
  \label{eq:defect-q}
  \delta_f^{(q)} \;:=\; \inf\bigl\{ t : P_{\mathrm{nom}}\bigl(\delta_f(x,u) \le t\bigr) \ge q \bigr\},
  \qquad 0 < q \le 1 .
\end{equation}
Then $\delta_f^{(q)} \le \delta_f^{(1)} \le \delta_f$ for all $q$;
$\delta_f^{(1)}$ is the essential supremum under $P_{\mathrm{nom}}$, and
$\delta_f$ is finite for any continuous $\hat f$ on the compact
$\mathcal{K}$. The sup version is not computable for a network; the
experiments report $\delta_f^{(0.95)}$ estimated on held-out nominal
validation samples (Sprint~6, Block~Q), and the corollaries below are stated
in both currencies. Equivariant models have $\delta_f = 0$ by
construction; for the mirror-shared DeLaN of \S\ref{sec:p2-learning} the
defect is exactly zero even in floating point (Remark~\ref{rem:exact-fp}).
\end{definition}

\begin{definition}[Residual and residual data element]
\label{def:residual-element}
The residual is $r(t) := y(t) - \hat f(x(t), u(t)) \in Y_r$. Its per-cycle
element is the phase-registered residual profile
$R_k : S^1 \to Y_r$, $R_k(\theta) = r(t_k(\theta))$ --- in implementation an
element of $\R^{d_r N}$ ($d_r$ residual rows on the $N$-point grid) ---
obtained from $Z_k$ by the \emph{residual map} $\mathfrak{r}$,
$(\mathfrak{r} Z)(\theta) := y(\theta) - \hat f(x(\theta), u(\theta))$,
applied pointwise in phase. $\symgrp$ acts on residual elements as on data
elements,
\begin{equation}
  \label{eq:res-action}
  (\rho_R(\sigma) R)(\theta) \;=\; \repR(g)\, R(\theta + \Delta\theta),
  \qquad \sigma = (g,\Delta\theta) \in \symgrp,
\end{equation}
so that for the trot $\rho_R := \rho_R(\sigma_*)$ is again an orthogonal
involution and $R = \Pi^{+}R + \Pi^{-}R$ with $\Pi^{\pm} = \tfrac12(I \pm
\rho_R)$, exactly as in \eqref{eq:rho-involution}. Wrapped and unfolded
versions of the element are defined as for $Z_k$ (Remark~\ref{rem:wrap}) ---
$\mathfrak{r}$ acts pointwise in phase and therefore commutes with both
constructions.
\end{definition}

\begin{remark}[Causal residual generators]
\label{rem:causal}
The analytic momentum observer is not a pointwise map: $r = K_O(p - \hat p)$
is a causal first-order filter of the data with cut-off $K_O$ (10\,Hz in
Sprint~4, i.e.\ a memory of $\approx 16$\,ms against a $500$\,ms cycle).
Everything below applies to a causal, time-invariant residual generator
$\mathfrak{r}$ that is $G$-equivariant, with two bookkeeping caveats: the
half-period shift commutes with a time-invariant filter, so
$\rho_R(\sigma)\mathfrak{r} = \mathfrak{r}\rho(\sigma)$ holds up to the
filter transient (negligible after a few time constants), and cross-cycle
independence of $(R_k)$ is inherited only up to that memory --- the same
exchangeability caveat that the raw data already carry
(Remark~\ref{rem:exch}). Under Assumption~\ref{as:dyn} the observer built
from the true model is $G$-equivariant: the equivariance of $M, C, g$ and of
the contact Jacobians is the statement of A1 for the rigid-body model
\cite{ordonez2025morphological}, and the observer is a fixed function of
them.
\end{remark}

\begin{remark}[The model-free channel is the special case $\hat f \equiv 0$]
\label{rem:model-free}
With $\hat f \equiv 0$ the residual is the torque block of the raw data,
$r = y$, and $R_k$ is the corresponding block of $Z_k$; the residual map is
a coordinate projection, trivially equivariant. Every statement of Part~1
about $Z_k$ is therefore the $\hat f \equiv 0$ instance of the statements
below --- the model-free $R^{-}$ channel of Sprints~1--2 is the degenerate
member of the family, and ``adding a model'' means replacing the projection
by $y - \hat f$.
\end{remark}

\subsection{Proposition N3-1: the residual inherits \texorpdfstring{$H_0$}{H0}}
\label{sec:p2-n31}

\begin{proposition}[N3-1, residual inherits $H_0$]
\label{prop:residual-h0}
Suppose the residual generator $\mathfrak{r}$ is $\symgrp$-equivariant on
cycle elements,
\begin{equation}
  \label{eq:r-equiv}
  \mathfrak{r}\bigl(\rho(\sigma) Z\bigr) \;=\; \rho_R(\sigma)\, \mathfrak{r}(Z)
  \qquad \forall \sigma \in \symgrp,
\end{equation}
which holds whenever $\hat f$ is $G$-equivariant \eqref{eq:equiv-model}
(pointwise generators), or is the observer of Remark~\ref{rem:causal} under
Assumption~\ref{as:dyn}. Then $H_0$ for the data element implies $H_0$ for
the residual element:
\begin{equation}
  \label{eq:h0-residual}
  \Law\bigl(\rho(\sigma) Z_k\bigr) = \Law(Z_k)\ \ \forall \sigma
  \quad\Longrightarrow\quad
  \Law\bigl(\rho_R(\sigma) R_k\bigr) = \Law(R_k)\ \ \forall \sigma \in \symgrp .
\end{equation}
Consequently Theorem~\ref{thm:n1-1} (exactness of the flip test) and
Theorem~\ref{thm:n1-2} (isotypic detectability) hold verbatim for the
sequence $(R_k)$ in place of $(Z_k)$: the flip test on residual elements is
exact under $H_0$, and a fault is visible to it iff its residual signature
has a nonzero sign component $\Pi^{-}\Delta_r \neq 0$.
\end{proposition}

\begin{proof}
Equivariance of the residual map: for a pointwise generator,
$(\mathfrak{r}\rho(\sigma)Z)(\theta) = \repY(g) y(\theta+\Delta\theta) -
\hat f(\repX(g)x(\theta+\Delta\theta), \repU(g)u(\theta+\Delta\theta)) =
\repR(g)\bigl[y - \hat f(x,u)\bigr](\theta+\Delta\theta) =
(\rho_R(\sigma)\mathfrak{r}Z)(\theta)$ by \eqref{eq:equiv-model} (the
$y$-block of $\repW$ restricted to $Y_r$ is $\repR$); compositions of
equivariant maps are equivariant. Distributional invariance passes through
an equivariant map: $\Law(\rho_R(\sigma)R_k) = \Law(\rho_R(\sigma)\mathfrak{r}(Z_k))
= \Law(\mathfrak{r}(\rho(\sigma)Z_k)) = \push{\mathfrak{r}}\Law(\rho(\sigma)Z_k)
= \push{\mathfrak{r}}\Law(Z_k) = \Law(R_k)$. Theorems~\ref{thm:n1-1}
and~\ref{thm:n1-2} assume only invariance of each element and independence
across elements; a measurable function of independent elements is
independent, so both hypotheses transfer.
\end{proof}

\begin{corollary}[Unified two-channel structure]
\label{cor:two-channel}
Under the hypotheses of Proposition~\ref{prop:residual-h0}, decompose
$R_k = \Pi^{+}R_k + \Pi^{-}R_k$. The $R^{-}$ statistic is a function of
$\Pi^{-}R_k$ only ($D_k = R_k - \rho_R R_k = 2\Pi^{-}R_k$), the magnitude
channel $R^{+}$ is a mirror-invariant function of $R_k$ whose informative
part under bilateral faults is $\Pi^{+}R_k$ (Theorem~\ref{thm:n1-2}(i)), and
$\|R_k\|^2 = \|\Pi^{+}R_k\|^2 + \|\Pi^{-}R_k\|^2$. The two channels are the
two isotypic components of \emph{one} residual read by two rules ---
symmetry breaking on $\Pi^{-}r$, magnitude on $\Pi^{+}r$ --- not two
statistics on two signals.
\end{corollary}

\begin{proof}
$\Pi^\pm$ are the orthogonal projectors of \eqref{eq:rho-involution} for
$\rho_R$; the identities are algebra. That $R^{-}$ ignores $\Pi^{+}R$ is
Theorem~\ref{thm:n1-2}(i) applied to $R_k$.
\end{proof}

\begin{remark}[Correction to the Part~0/1 asymmetry]
Part~0 (Example~\ref{ex:bilateral}) and the Sprint~2/4 pipelines paired an
$R^{-}$ test on the \emph{raw} signal with an $R^{+}$ score on the
\emph{residual}. Corollary~\ref{cor:two-channel} shows the natural object is
the residual for both: $R^{-}$ on $\Pi^{-}r$ inherits exactness from
Proposition~\ref{prop:residual-h0} and, by Proposition~\ref{prop:power-gain},
generally more power than on $y$; the raw-signal $R^{-}$ remains available
as the $\hat f \equiv 0$ member (Remark~\ref{rem:model-free}), and it is the
only member that needs no model at all --- which is why it stays the
deployment fallback when no nominal model is trusted.
\end{remark}

\begin{corollary}[A non-equivariant model contaminates $H_0$]
\label{cor:contamination}
Let $\hat f$ have defect $\delta_f > 0$ and let $\hat f^{\mathrm{sym}} :=
|G|^{-1}\sum_{g\in G} \repR(g)^{-1}\hat f(\repX(g)\cdot,\repU(g)\cdot)$ be its
group average, which is $G$-equivariant. Then $\hat f = \hat f^{\mathrm{sym}}
+ \Delta\hat f$ with $\|\Delta \hat f(x,u)\| \le \tfrac{|G|-1}{|G|}\,\delta_f(x,u)$
($=\tfrac12\delta_f(x,u)$ for $\mathbb{C}_2$), and the residual splits as
$R_k = R_k^{\mathrm{sym}} - \Delta F(Z_k)$, where $R^{\mathrm{sym}}_k :=
\mathfrak{r}^{\mathrm{sym}}(Z_k)$ satisfies $H_0$ exactly under the
hypotheses of Proposition~\ref{prop:residual-h0} and $\Delta F(Z)(\theta) =
\Delta\hat f(x(\theta),u(\theta))$. For the actual residual,
\begin{equation}
  \label{eq:eps-model}
  \begin{aligned}
  \dtv\bigl(\Law(\rho_R R_k), \Law(R_k)\bigr) \;&\le\;
  \epsbar{model} + \epsbar{tot},\\
  \epsbar{model} \;&:=\; \sup_{\sigma\in\symgrp}
  \dtv\bigl(\Law(\mathfrak{r}(\rho(\sigma)Z_k)),\ \Law(\rho_R(\sigma)\mathfrak{r}(Z_k))\bigr),
  \end{aligned}
\end{equation}
so $\epsbar{model}$ enters Lemma~\ref{lem:eps} as a fifth additive term and
the level bound becomes $\mathbb{E}[\varphi] \le \alpha + \tfrac12
(\epsbar{tot} + \epsbar{model})$ for the trot. It is controlled by the
defect: writing $e(Z) := \mathfrak{r}(\rho Z) - \rho_R\mathfrak{r}(Z)$, one has
$\|e(Z)(\theta)\| = \delta_f(x(\theta),u(\theta))$ pointwise, and
\begin{equation}
  \label{eq:c-delta}
  \epsbar{model} \;\le\; c\,\delta_f,
  \qquad
  c \;:=\; \sup_{0 < \|v\|_{\infty} \le \delta_f}
  \frac{\dtv\bigl(\Law(R_k^{\mathrm{sym}} + v),\ \Law(R_k^{\mathrm{sym}})\bigr)}{\|v\|_{\infty}} ,
\end{equation}
the total-variation modulus of the nominal residual law under bounded
deterministic shifts (norm $\|v\|_\infty := \max_\theta \|v(\theta)\|$). If
the residual carries an additive noise component $\zeta \sim
\mathcal{N}(0, s^2 I_{d_r N})$ independent of the rest (e.g.\ torque
measurement or actuator noise passed through the residual map), then
$c \le \sqrt{N}/(s\sqrt{2\pi})$ and $\epsbar{model} \le \sqrt{N}\,\delta_f/(s\sqrt{2\pi})$.
In the quantile currency, with $\delta_f^{(q)}$ of \eqref{eq:defect-q},
\begin{equation}
  \label{eq:c-delta-q}
  \epsbar{model} \;\le\; N(1-q) + c\,\delta_f^{(q)} ,
\end{equation}
and if the quantile is taken over cycle elements ($\max_\theta$ of the
pointwise defect) the first term is $(1-q)$. Equivariant models have
$\epsbar{model} = 0$ by construction.
\end{corollary}

\begin{proof}
Group averaging: by re-indexing the sum,
\[
  \repR(h)^{-1}\hat f^{\mathrm{sym}}(\repX(h)\cdot,\repU(h)\cdot)
  = |G|^{-1}\sum_g \repR(hg)^{-1}\hat f(\repX(hg)\cdot,\repU(hg)\cdot) = \hat f^{\mathrm{sym}},
\]
and $\Delta\hat f = |G|^{-1}\sum_{g \neq e}[\hat f -
\repR(g)^{-1}\hat f(\repX(g)\cdot,\repU(g)\cdot)]$ has $|G|-1$ terms each
bounded by $\delta_f(x,u)$ (orthogonality of $\repR(g)$). The split of $R_k$
is linearity of $\mathfrak{r}$ in $\hat f$, and $R^{\mathrm{sym}}_k$
inherits $H_0$ by Proposition~\ref{prop:residual-h0}. For
\eqref{eq:eps-model}: $\dtv(\Law(\rho_R R), \Law(R)) \le \dtv(\Law(\rho_R
\mathfrak{r}Z), \Law(\mathfrak{r}\rho Z)) + \dtv(\Law(\mathfrak{r}\rho Z),
\Law(\mathfrak{r} Z))$; the second term is at most $\dtv(\Law(\rho Z),
\Law(Z)) \le \epsbar{tot}$ by the data-processing inequality and
Lemma~\ref{lem:eps}, the first is $\epsbar{model}$ by definition. The level
statement is Lemma~\ref{lem:eps}(ii) with the enlarged budget. For
\eqref{eq:c-delta}: $\mathfrak{r}(\rho Z) = \rho_R \mathfrak{r}(Z) + e(Z)$
with $e(Z)(\theta) = \repR\hat f(x,u) - \hat f(\repX x,\repU u)$ evaluated at
$\theta + \Delta\theta$, of norm $\delta_f(x,u) \le \delta_f$; conditioning
on the value of $e$ (a bounded, $Z$-measurable shift) and using
$\Law(\rho_R\mathfrak{r} Z) = \Law(\rho_R R^{\mathrm{sym}} - \rho_R\Delta F(Z))$
together with the definition of $c$ as the worst-case modulus over shifts of
sup-norm at most $\delta_f$ gives $\epsbar{model} \le c\,\delta_f$ (the
shift $-\rho_R \Delta F(Z)$ has sup-norm $\le \tfrac12\delta_f$ and $e$ at
most $\delta_f$; the modulus is monotone in the shift). The Gaussian
constant: for $\zeta \sim \mathcal{N}(0,s^2 I_n)$ and any $v \in \R^n$,
$\dtv(\mathcal{N}(v,s^2 I), \mathcal{N}(0,s^2 I)) = 2\Phi(\|v\|_2/2s) - 1 \le
\|v\|_2/(s\sqrt{2\pi})$, and $\|v\|_2 \le \sqrt{N}\,\|v\|_\infty$ for a
profile of $N$ phase samples; conditioning on everything but $\zeta$ and
averaging (convexity of $\dtv$) gives the bound. For \eqref{eq:c-delta-q}:
the event $\{\exists\theta : \delta_f(x(\theta),u(\theta)) > \delta_f^{(q)}\}$
has probability at most $N(1-q)$ by a union bound over the $N$ grid points
(each marginal is nominal); on its complement the previous argument applies
with $\delta_f^{(q)}$ in place of $\delta_f$, and total variation splits over
the two events.
\end{proof}

\paragraph{Falsification condition (N3-1).}
Assumption~\ref{as:dyn} failing (the true dynamics not $G$-equivariant, e.g.\
a manufacturing asymmetry above $\tau_{\mathrm{dyn}}$) breaks $H_0$ for
$Z_k$ and hence for $R_k$ --- with an \emph{equivariant} $\hat f$ the
residual then carries the plant asymmetry, which no model choice can remove
(the observer built from the symmetrized model is equivariant, the plant is
not). The degradation path is Part~0's: $H_0'$ on the residual (calibration
of the nominal asymmetry level $\nu_0$; for the residual this is the
centring of Lemma~\ref{lem:centring}). Empirical anchor: e13b (size of the
residual flip test under $H_0$ against $\delta_f^{(0.95)}$ of the model:
plain DeLaN ladder vs equivariant DeLaN vs analytic residual;
Section~\ref{sec:p2-anchors}). What would falsify the corollary's mechanism
is a plain model with large measured $\delta_f^{(0.95)}$ whose residual flip
test keeps nominal size at all $K$; what would falsify
Proposition~\ref{prop:residual-h0} is an equivariant residual whose flip
test inflates although the raw-signal test does not.

\subsection{Proposition N3-2: power gain of the residual channel}
\label{sec:p2-n32}

\paragraph{Setting.}
Write the torque-block of the data as $y = f^{*}(x,u) + \zeta$ with $f^{*}$
the true nominal map and $\zeta$ the fluctuation not explained by the state
(measurement and actuator noise, unmodelled contact softness), and the model
as $\hat f = f^{*} + m$ with model error $m$. A fault $\varphi$ shifts the
nominal mean trajectory of the data element by $\Delta_y(\varphi)$ (torque
block) and of the residual element by $\Delta_r(\varphi)$; both are profiles
in $\R^{d_r N}$ and both decompose isotypically. Two channel-wise scalars
summarize what a standardized energy statistic sees:
\begin{equation}
  \label{eq:snr}
  \mathrm{SNR}_y \;:=\; \sum_{c} \frac{(\Pi^{-}\Delta_y)_c^{2}}{\operatorname{Var}(\Pi^{-}y)_c},
  \qquad
  \mathrm{SNR}_r \;:=\; \sum_{c} \frac{(\Pi^{-}\Delta_r)_c^{2}}{\operatorname{Var}(\Pi^{-}r)_c},
\end{equation}
sums over the residual rows $c$ (phase samples included in the index), where
$\operatorname{Var}(\Pi^{-}\cdot)_c$ is the nominal per-channel variance of
the sign component --- the per-channel standardization of the deployed
statistic (Section~\ref{sec:audit}: pooled flip-invariant scale of the
calibration cycles).

\begin{proposition}[N3-2, power gain of the residual channel]
\label{prop:power-gain}
Consider location alternatives as in Theorem~\ref{thm:n1-2}(ii) with
Gaussian fluctuations, the standardized paired-energy statistic, and $K$
cycles. Then:
\begin{enumerate}
\item[(a)] \emph{(fault-model commutation)} If the fault acts \emph{after}
the command in the actuation chain --- gain $\kappa$, bias $b$, or a
friction change on joint $j$, with $y$ built from the \emph{commanded}
torque --- and the closed-loop state trajectory is unchanged to first order,
then
\begin{equation}
  \label{eq:delta-r}
  \Delta_r \;=\; \Delta_y \;-\; \bigl[\hat f(x_{\varphi},u_{\varphi}) - \hat f(x_{\mathrm{nom}},u_{\mathrm{nom}})\bigr]
  \;=\; \Delta_y \quad\text{(response term zero)},
\end{equation}
and the residual signature is the fault's own torque footprint:
$\Delta_r = (1-\kappa)\,\bar\tau_j(\theta)\,e_j$, $\Delta_r = -b\,e_j$,
$\Delta_r = -\Delta\mu\,\operatorname{sign}\bar{\dot q}_j(\theta)\,e_j$
respectively (the rank-1 dictionary of the N3 certificate, Sprint~4).
\item[(b)] \emph{(variance removal)} If the model explains part of the
nominal antisymmetric fluctuation, $\operatorname{Var}(\Pi^{-}r)_c \le
\operatorname{Var}(\Pi^{-}y)_c$ for every row $c$ --- which holds when
$m + \zeta$ has less antisymmetric variance than $f^{*}(x,u) + \zeta$, i.e.\
when the state-driven part of the torque fluctuation dominates the model
error --- then, under (a),
\begin{equation}
  \label{eq:snr-ratio}
  \frac{\mathrm{SNR}_r}{\mathrm{SNR}_y}
  \;=\;
  \frac{\sum_c (\Pi^-\Delta_y)_c^2 / \operatorname{Var}(\Pi^-r)_c}
       {\sum_c (\Pi^-\Delta_y)_c^2 / \operatorname{Var}(\Pi^-y)_c}
  \;\ge\; \min_{c : (\Pi^-\Delta_y)_c \neq 0} \frac{\operatorname{Var}(\Pi^-y)_c}{\operatorname{Var}(\Pi^-r)_c}
  \;\ge\; 1 ,
\end{equation}
and since the power of the standardized paired-energy flip test is
increasing in the noncentrality $K \cdot \mathrm{SNR}$
(Theorem~\ref{thm:n1-2}(ii)), the residual channel is at least as powerful
as the raw torque channel at every $K$, with the gain given by the share of
nominal antisymmetric variance that $\hat f$ explains.
\item[(c)] \emph{(general faults, conditional)} Without (a) the response term
in \eqref{eq:delta-r} is present, $\Pi^-\Delta_r = \Pi^-\Delta_y -
\Pi^-\Delta_{\hat f}$, and \eqref{eq:snr-ratio} holds with
$(\Pi^-\Delta_y)_c^2$ replaced by $(\Pi^-\Delta_r)_c^2$ in the numerator: the
residual channel wins iff the response term does not cancel the fault
footprint by more than the variance factor,
$\sum_c (\Pi^-\Delta_r)_c^2/\operatorname{Var}(\Pi^-r)_c \ge
\sum_c (\Pi^-\Delta_y)_c^2/\operatorname{Var}(\Pi^-y)_c$.
\end{enumerate}
\end{proposition}

\begin{proof}
(a) With $y = \tau_{\mathrm{cmd}} + J^\top f_{\mathrm{c}}$ and the physics
$\tau_{\mathrm{app}} + J^\top f_{\mathrm{c}} = f^{*}(x,u) + \zeta$, the
residual is $r = \tau_{\mathrm{cmd}} - \tau_{\mathrm{app}} - m(x,u) + \zeta$;
gain and bias act on $\tau_{\mathrm{app}} = \kappa\tau_{\mathrm{cmd}} + b$
and a friction change adds $-\Delta\mu\operatorname{sign}\dot q_j$ to the
torque balance, so their footprint enters $r$ additively and $\Delta_r$ is
their mean profile; $\Delta_y$ is the same footprint (the torque block of
$Z$ shifts by the same amount when the state is unchanged) and
$\hat f(x_\varphi,u_\varphi) = \hat f(x_{\mathrm{nom}},u_{\mathrm{nom}})$
by the first-order assumption. (b) is a termwise comparison of the two sums
in \eqref{eq:snr}: each term of $\mathrm{SNR}_r$ is the corresponding term
of $\mathrm{SNR}_y$ times $\operatorname{Var}(\Pi^-y)_c/\operatorname{Var}(\Pi^-r)_c
\ge 1$; monotonicity of power in the noncentrality is the Gaussian
noncentral-$\chi^2$ (or its permutation analogue) monotonicity in the
location model. (c) is the same computation without the identification
$\Delta_r = \Delta_y$.
\end{proof}

\begin{remark}[When the residual is \emph{not} better]
\label{rem:residual-weaker}
Three mechanisms make the residual channel weaker than the raw channel, and
each is a case of (c): (i) \emph{closed-loop response} --- a gain fault
changes the tracking error, the controller changes $\tau_{\mathrm{cmd}}$ and
the trajectory, and $\hat f$ responds; if $\Pi^-\Delta_{\hat f}$ is aligned
with $\Pi^-\Delta_y$ the footprint is partly cancelled (the residual sees the
\emph{unexplained} torque, and a compensating controller moves part of the
fault into explained torque); (ii) \emph{faults outside the torque balance}
--- an encoder bias shifts the raw $q$-channel by $b$ directly, while it
reaches the residual only through $\hat f(q+b) - \hat f(q) \approx
\partial_q\hat f\, b$, typically much smaller after standardization; (iii)
\emph{channel count} --- the raw element carries $q,\dot q$, IMU and contact
rows that respond to the fault kinematically, the residual only its $d_r$
torque rows, so the system-level comparison of the two channels (all raw
channels vs all residual rows, as run in e13a) is not the row-wise
comparison of \eqref{eq:snr-ratio}. Any grid cell in which the residual
version is weaker is to be reported as an instance of this remark with the
mechanism identified, not tuned away.
\end{remark}

\paragraph{Falsification condition and empirical anchor (N3-2).}
Anchor: e13a --- residual $R^{-}$ (analytic and equivariant DeLaN) versus
raw-signal $R^{-}$ on the low-SNR grid (gain $1-\kappa \in \{0.01, 0.02,
0.05, 0.1\}$, bias $\{0.05, 0.1, 0.2, 0.5\}$\,N\,m, friction $\times\{1.1,
1.2, 1.5\}$ on LF-HFE/KFE, $R = 50$), with the nominal variance
decomposition $\operatorname{Var}(\Pi^-y)_c$ vs $\operatorname{Var}(\Pi^-r)_c$
row by row and the Mahalanobis magnitude reference
(Section~\ref{sec:p2-anchors}). The proposition is falsified in a cell if
the variance condition (b) holds, the fault is of type (a), and the residual
channel is nevertheless less powerful; a residual loss with the variance
condition violated, or for a fault of type (c), is the conditional case of
Remark~\ref{rem:residual-weaker}.

\subsection{Where the learned model sits (N3 updated)}
\label{sec:p2-learning}

The DK isolability certificate of Sprint~4 (\texttt{geofdi.isolation.n3};
\cite{yu2015daviskahan}) reads the residual through a rank-1 signature
dictionary and a perturbation constant $\beta_{\mathrm{op}}^2$, the top
eigenvalue of the second moment of the \emph{calibration-centred} nominal
residual profiles. Corollary~\ref{cor:two-channel} places this constant: it
is measured on the fluctuating part of $R_k$ --- both isotypic components,
since the certificate reads profiles, not tests --- and it is at once the
width of the $R^{+}$ conformal calibration (the nominal residual magnitude
the conformal $p$-value is compared with) and the perturbation term of the
certificate. Two things must be said about it: why systematic model error
does not enter it (Sprint~4, Finding~4), and how the learned model is made
equivariant so that Proposition~\ref{prop:residual-h0} applies to it.

\begin{lemma}[Calibration centring removes systematic model error]
\label{lem:centring}
Let the residual element decompose as $R_k = \bar m + \Xi_k$ where
$\bar m \in \R^{d_r N}$ is a deterministic gait-locked profile (systematic
model error and observer floor: a function of phase through the nominal
mean trajectory) and $\Xi_k$ are i.i.d.\ zero-mean fluctuations with
covariance $\Sigma_\Xi$; under a fault of magnitude $m$ with signature
$s_c$, $R_k = \bar m + m\,s_c + \Xi_k$. Let $\hat{\bar m} :=
K_{\mathrm{cal}}^{-1}\sum_{k \le K_{\mathrm{cal}}} R_k$ be the calibration
mean. Then:
\begin{enumerate}
\item[(i)] under nominal operation the centred profiles $\tilde R_k := R_k -
\hat{\bar m}$ have $\mathbb{E}\tilde R_k = 0$ and
$\operatorname{Var}(\tilde R_k) = \Sigma_\Xi (1 + K_{\mathrm{cal}}^{-1})$, so
$\beta_{\mathrm{op}}^2 = \lambda_{\max}(\Sigma_\Xi)(1+K_{\mathrm{cal}}^{-1})$
does not depend on $\bar m$;
\item[(ii)] under the fault the population second moment of $\tilde R_k$ is
$m^2 s_c s_c^\top + \Sigma_\Xi(1+K_{\mathrm{cal}}^{-1})$ --- the certificate's
setting with $\beta_{\mathrm{op}}^2$ as in (i) and the fault signature
intact; for the $R^{+}$ conformal channel and the isolation classifier,
which compare monitored profiles with the calibration distribution of the
same quantity, the estimation error of $\hat{\bar m}$ is part of both sides
and is harmless;
\item[(iii)] \emph{(naive centring does not rescue the $R^{-}$ flip test)}
$\Pi^{-}\tilde R_k = \Pi^{-}\Xi_k - o$ with the \emph{common} offset
$o := \Pi^{-}(\hat{\bar m} - \bar m)$ of covariance $\Sigma^{-}_\Xi/K_{\mathrm{cal}}$,
so the systematic asymmetry $\Pi^{-}\bar m$ of a non-equivariant model is
removed in the mean; but the flip test on $\tilde R_k$ is not exact: for the
paired-energy statistic of $K$ monitored cycles the observed value exceeds
its flip-null counterpart in expectation by the factor
$(K_{\mathrm{cal}} + K)/(K_{\mathrm{cal}} + 1) \approx 1 + K/K_{\mathrm{cal}}$,
while the null statistic aggregates $d_r N$ coordinates and has relative
spread $O((d_r N)^{-1/2})$; the size is nominal only if
$K\sqrt{d_r N} \ll K_{\mathrm{cal}}$, which no practical calibration window
satisfies ($K = K_{\mathrm{cal}} = 60$, $d_r N = 768$: a factor-2 shift
against a $5\,\%$ spread, size $\approx 1$). Second-order asymmetries
(mirror-unequal fluctuation laws of a non-equivariant model) survive
centring in any case;
\item[(iv)] \emph{(the $H_0'$ construction for $R^{-}$)} difference each
monitored cycle with an independent calibration cycle,
$X_k := R^{\mathrm{mon}}_k - R^{\mathrm{cal}}_k$, $k \le K \le K_{\mathrm{cal}}$.
Under $H_0'$ (calibration and monitoring have the same law, whatever the
model's defect) the paired differences $D_k = X_k - \rho_R X_k = A_k - B_k$
with $A_k := R^{\mathrm{mon}}_k - \rho_R R^{\mathrm{mon}}_k$ and
$B_k := R^{\mathrm{cal}}_k - \rho_R R^{\mathrm{cal}}_k$ i.i.d., are exactly
sign-symmetric, $D_k \eqdist -D_k$, and independent across $k$; hence the
flip test with any statistic of the tuple $(D_1,\dots,D_K)$ --- the
paired-energy statistic in particular --- is exact under $H_0'$ by the argument of
Theorem~\ref{thm:n1-1} (the flip group acts as $D_k \mapsto \varepsilon_k D_k$
and $\Law(\varepsilon D) = \Law(D)$). The price is one calibration cycle per
monitored cycle and a doubled fluctuation variance (halved SNR); a fault
appears as a mean of $D_k$ equal to the change of the asymmetry between
calibration and monitoring, which is Part~0's $H_0'$ reading. Statistics that
flip the elements themselves (energy distance between $\{X_k\}$ and
$\{\rho_R X_k\}$) are not covered, since $\Law(X_k)$ need not be
$\rho_R$-invariant.
\end{enumerate}
\end{lemma}

\begin{proof}
(i)--(ii): $\hat{\bar m} = \bar m + \bar\Xi_{\mathrm{cal}}$ with
$\bar\Xi_{\mathrm{cal}}$ independent of the monitored $\Xi_k$ and of
covariance $\Sigma_\Xi/K_{\mathrm{cal}}$; the systematic profile cancels
identically and the variances add. (iii): with $\bar D := K^{-1}\sum_k
(\Pi^-\Xi_k - \rho_R\Pi^-\Xi_k)$ of per-coordinate variance $2\sigma^2/K$ and
$o' := o - \rho_R o$ of per-coordinate variance $2\sigma^2/K_{\mathrm{cal}}$,
$\mathbb{E}\|\bar D - o'\|^2 = d_r N\,(2\sigma^2/K + 2\sigma^2/K_{\mathrm{cal}})$,
whereas for a uniform flip pattern $\varepsilon$ with $\bar\varepsilon :=
K^{-1}\sum_k \varepsilon_k$ ($\mathbb{E}\bar\varepsilon^2 = 1/K$) the flipped
statistic $\|K^{-1}\sum_k \varepsilon_k D_k - \bar\varepsilon o'\|^2$ has
expectation $d_r N\,(2\sigma^2/K + 2\sigma^2/(K K_{\mathrm{cal}}))$; the
ratio is $(K_{\mathrm{cal}}+K)/(K_{\mathrm{cal}}+1)$, and a sum of $d_r N$
weakly dependent squared terms has relative standard deviation of order
$(d_r N)^{-1/2}$. (iv): $A_k \eqdist B_k$ and independence give
$A_k - B_k \eqdist B_k - A_k = -(A_k - B_k)$; independence across $k$ is
inherited from the cycles; the rest is Theorem~\ref{thm:n1-1} with the group
$\{\pm 1\}^K$ acting on $D$.
\end{proof}

(i)--(ii) are the mechanism behind Sprint~4, Finding~4 (degrading a DeLaN
by sample size raised $\hat\beta$ five-fold but $\beta_{\mathrm{op}}$ only by
$1.7\times$: the systematic part was calibrated away). (iii)--(iv) say how
Part~0's $H_0'$ route applies to the $R^{-}$ channel: not by subtracting an
estimated profile and re-running the exact test, but by differencing against
calibration cycles --- exact for any model, at half the SNR. Empirical anchor
(e13b, Section~\ref{sec:p2-anchors}): naive centring gives size $\approx 1$
for every model, the differenced test is in band for every model including
the plain DeLaN with $\delta_f^{(0.95)} = 15$\,N\,m. Both limits are why the
equivariant model of the next proposition is preferable for the $R^{-}$
channel: with $\delta_f = 0$ nothing needs to be calibrated away and the
full-SNR exact test of Proposition~\ref{prop:residual-h0} applies.

\begin{proposition}[Equivariant DeLaN by mirror weight sharing]
\label{prop:equiv-delan}
Let $\hat f_0 : \R^{3}\times\R^{3}\times\R^{3}\times\R^{3} \to \R^{3}$,
$(q,\dot q,\ddot q, a) \mapsto \tau$, be a leg Lagrangian network
\cite{lutter2019delan} for the template leg (LF, resp.\ LH): with
$M_0(q) = L(q)L(q)^\top + \epsilon I$ ($L$ lower triangular with
softplus-positive diagonal), potential $V_0(q;a)$ parametrized by the trunk
specific force $a$, and fixed passive terms $b\dot q + f\tanh(\dot q/c)$
($b, f$ diagonal, $c > 0$),
\begin{equation}
  \label{eq:delan-template}
  \hat f_0(q,\dot q,\ddot q,a) \;=\; M_0(q)\ddot q + C_0(q,\dot q)\dot q +
  \nabla_q V_0(q;a) + b\dot q + f\tanh(\dot q/c),
\end{equation}
$C_0$ the Christoffel matrix of $M_0$. Let $S = \operatorname{diag}(s_{\mathrm{HAA}},
s_{\mathrm{HFE}}, s_{\mathrm{KFE}})$ be the joint-sign matrix of the sagittal
reflection (Table~\ref{tab:joint-signs}; $S = \operatorname{diag}(-1,1,1)$
in the uniform-axis convention) and $E = \operatorname{diag}(1,-1,1)$ its
action on the specific force. Define the mirrored leg's model by
\begin{equation}
  \label{eq:mirror-template}
  \hat f_{\sigma}(q,\dot q,\ddot q,a) \;:=\; S\, \hat f_0\bigl(Sq,\; S\dot q,\; S\ddot q,\; Ea\bigr),
\end{equation}
and the four-leg model by $\hat F := (\hat f_{\mathrm{LF}}, \hat f_{\mathrm{RF}},
\hat f_{\mathrm{LH}}, \hat f_{\mathrm{RH}}) = (\hat f_{0}^{F}, (\hat f_0^{F})_\sigma,
\hat f_0^{H}, (\hat f_0^{H})_\sigma)$ with one template per front/hind pair.
Then:
\begin{enumerate}
\item[(i)] $\hat F$ is exactly $\mathbb{C}_2$-equivariant: $\hat F(\repX(\gs)z)
= \repR(\gs)\hat F(z)$ for all $z$, so $\delta_f = 0$.
\item[(ii)] $\hat f_\sigma$ is the inverse dynamics of the Lagrangian
$L_\sigma(q,\dot q;a) := L_0(Sq, S\dot q; Ea)$: its mass matrix
$M_\sigma(q) = S M_0(Sq) S$ is symmetric positive definite, its potential is
$V_\sigma(q;a) = V_0(Sq;Ea)$, its Coriolis matrix is the Christoffel matrix
of $M_\sigma$, and the passive terms are unchanged. The mirrored leg
therefore has the same physical structure as the template (positive
definiteness, Christoffel consistency $\dot M_\sigma = C_\sigma + C_\sigma^\top$,
conservative potential) --- inherited, not re-learned.
\item[(iii)] With two templates the construction enforces exactly the
morphological symmetry $G = \mathbb{C}_2$ of a robot whose front and hind
pairs differ (M1, Go2); a single template for all four legs
(\texttt{n\_templates}$=1$) additionally imposes a front--hind
identification that is \emph{not} a symmetry of these robots and is used
only as an ablation.
\end{enumerate}
\end{proposition}

\begin{proof}
(i) $\repX(\gs)$ sends the LF data $(q,\dot q,\ddot q)$ to the RF slot with
the signs $S$ and $a \mapsto Ea$ (Section~\ref{sec:reps}), and $\repR(\gs)$
sends the LF torque to the RF slot with $S$. The RF component of
$\hat F(\repX(\gs)z)$ is $\hat f_\sigma(Sq_{\mathrm{LF}}, S\dot q_{\mathrm{LF}},
S\ddot q_{\mathrm{LF}}, Ea) = S\hat f_0(S^2 q_{\mathrm{LF}}, \dots, E^2 a) =
S\hat f_0(q_{\mathrm{LF}},\dots,a) = S\hat f_{\mathrm{LF}}(z)$, which is the RF
component of $\repR(\gs)\hat F(z)$; the LF component of $\hat F(\repX(\gs)z)$
is $\hat f_0(Sq_{\mathrm{RF}}, \dots, Ea) = S\cdot S\hat f_0(Sq_{\mathrm{RF}},\dots,Ea)
= S\hat f_\sigma(q_{\mathrm{RF}},\dots,a) = S \hat f_{\mathrm{RF}}(z)$, the LF
component of $\repR(\gs)\hat F(z)$; the hind pair is identical
($S^2 = E^2 = I$). (ii) With $q' := Sq$ (so $\dot q' = S\dot q$, $S = S^\top =
S^{-1}$), $L_\sigma(q,\dot q;a) = \tfrac12 \dot q'^\top M_0(q')\dot q' -
V_0(q';Ea) = \tfrac12\dot q^\top S M_0(Sq) S\,\dot q - V_0(Sq;Ea)$, so
$M_\sigma = S M_0(S\cdot) S$: a congruence of an SPD matrix by an
orthogonal matrix, hence SPD, and $V_\sigma(q;a) = V_0(Sq;Ea)$. By the chain
rule $\nabla_q V_\sigma = S(\nabla_{q'}V_0)(Sq;Ea)$ and the Euler--Lagrange
expression transforms as $\frac{d}{dt}\partial_{\dot q}L_\sigma -
\partial_q L_\sigma = S\bigl[\frac{d}{dt}\partial_{\dot q'}L_0 -
\partial_{q'}L_0\bigr](Sq,S\dot q,S\ddot q)$, i.e.\ the inverse dynamics of
$L_\sigma$ is $S$ times the inverse dynamics of $L_0$ at the mirrored
arguments; the Christoffel matrix of $M_\sigma$ is $C_\sigma(q,\dot q) =
S\,C_0(Sq,S\dot q)\,S$ (Christoffel symbols are built from first
derivatives of $M$, and $\partial_{q_i} = s_i\partial_{q'_i}$), so
$C_\sigma\dot q = S C_0(Sq,S\dot q)S\dot q$, matching \eqref{eq:mirror-template}
term by term. Passive terms: $S b S\dot q = b\dot q$ for diagonal $b$, and
$S f\tanh(S\dot q/c) = f\tanh(\dot q/c)$ because $\tanh$ is odd and $S$
diagonal with entries $\pm 1$. Since the DeLaN forward pass computes exactly
the Euler--Lagrange expression of its template
\cite{lutter2019delan}, \eqref{eq:mirror-template} \emph{is} the inverse
dynamics of $L_\sigma$. (iii) The reflection pairs LF$\leftrightarrow$RF and
LH$\leftrightarrow$RH and nothing else (Section~\ref{sec:m1}); a shared
template across the pairs would identify front and hind leg subsystems.
\end{proof}

\begin{remark}[Exactly zero, also in floating point]
\label{rem:exact-fp}
$S$ and $E$ are signed permutations, and multiplication by $\pm 1$ is exact
in floating point; the defect $\hat f_\sigma(\repX z) - S\hat f_0(z)$ is
computed by evaluating the \emph{same} network on bitwise-identical inputs
($S^2 q = q$ exactly), so $\delta_f$ measured on data is $0$ to the last
bit --- not $10^{-6}$ (Sprint~6 unit tests \texttt{tests/test\_delan\_equiv.py}
and the $\delta_f$ column of Block~Q). Weight sharing across a leg orbit is
the discrete-group instance of group-equivariant architectures
\cite{cohen2016group,weiler2019general,finzi2021practical} --- a group
convolution over the leg index --- and is how morphological symmetries are
imposed on learned dynamics in \cite{ordonez2025morphological}; here the
convolution has a single orbit per template and no learnable mixing, so
equivariance costs nothing and no basis of equivariant maps has to be
solved for.
\end{remark}

\begin{remark}[Why not a soft (PINN-style) penalty]
\label{rem:no-soft}
Physical consistency comes from the architecture (DeLaN: $M = LL^\top +
\epsilon I$ is positive definite for every input, the Coriolis term is the
Christoffel term of that $M$), symmetry consistency from weight sharing.
Neither is a loss term, and for a reason that matters to the theory: a
penalty $\lambda\,\mathbb{E}\|\hat f(\rho z) - \rho\hat f(z)\|^2$ makes the
defect small \emph{on the training distribution and in expectation}, but
gives no bound on $\delta_f$ (nor on $\delta_f^{(q)}$ off the training
support), so $\epsbar{model}$ in Corollary~\ref{cor:contamination} cannot be
written down; likewise a positivity penalty on $M$ can be violated at test
time, and Proposition~\ref{prop:equiv-delan}(ii) would become a statement
about training loss rather than about the model. The residual $R^{-}$ test
is a claim about the null \emph{law}; only hard structure gives it something
to inherit.
\end{remark}

\paragraph{Falsification condition and empirical anchor (equivariant DeLaN).}
The proposition is exact; what is empirical is whether the constraint costs
accuracy: the anchor is Block~Q --- validation RMSE of the equivariant
ladder (full, $50$k, $10$k, $2$k samples per leg) against the plain ladder,
and $\delta_f^{(0.95)}$ of both (Section~\ref{sec:p2-anchors}). A
validation RMSE more than $10\,\%$ worse than the plain model would indicate
that a shared template lacks capacity for the two legs of a pair (or that
the pair is not physically mirror-symmetric --- an A1 signal on the training
world), and would be reported as such.

\subsection{Updated \texorpdfstring{$\varepsilon$}{epsilon}-budget and the three-channel architecture}
\label{sec:p2-budget}

Table~\ref{tab:degradation2} extends the falsification--degradation map of
Table~\ref{tab:degradation} by the model row. Figure~\ref{fig:architecture}
draws the three readouts of the residual --- $R^{-}$ on $\Pi^{-}r$, the
joint residual rows (with $\Pi^{+}r$ as their magnitude reading and the
per-leg deviation for left/right attribution), and the six floating-base
rows of the analytic observer (payload and other base-level nuisances,
Sprint~4 Finding~2) --- and marks the one learned component. Everything
downstream of $\hat f$ has zero trainable parameters: the flip test, the
conformal calibration, the e-process/e-CUSUM layer, and the isolation rules
(pair--joint from the $R^{-}$ ranking, left/right from the joint rows,
payload from the base rows).

\begin{table}[ht]
\centering
\small
\begin{tabular}{@{}p{0.17\linewidth}p{0.36\linewidth}p{0.40\linewidth}@{}}
\toprule
term & falsification signal & degradation path \\
\midrule
$\epsbar{dyn}$, $\epsbar{ctrl}$, $\epsbar{env}$, $\epsbar{noise}$ &
as in Table~\ref{tab:degradation} & as in Table~\ref{tab:degradation};
all four pass into the residual unchanged (data-processing inequality) \\
$\epsbar{model}$ (new) &
$\delta_f^{(0.95)} > 0$ on mirrored validation data (Block~Q); residual
flip-test size above the band on nominal data while the raw test is in band
(e13b) &
(1) replace $\hat f$ by its equivariant version (weight sharing,
Proposition~\ref{prop:equiv-delan}): $\epsbar{model} = 0$;
(2) $H_0'$ on the residual by differencing against calibration cycles
(Lemma~\ref{lem:centring}(iv); exact for any model, half the SNR --- naive
mean-profile centring is \emph{not} valid for the flip test, (iii)); (3) budget it: level
$\alpha + \tfrac12(\epsbar{tot} + \epsbar{model})$ with
\eqref{eq:c-delta}--\eqref{eq:c-delta-q} \\
$\epsbar{chir}$ ($\nu_0$) & Gate~1b & unchanged: $H_0'$ on the residual
(centring is its calibration) \\
residual memory & causal generator with memory comparable to the cycle &
unfolded elements / block flips (Remark~\ref{rem:causal}) \\
\bottomrule
\end{tabular}
\caption{Falsification--degradation map extended by the model term of
Corollary~\ref{cor:contamination}. The four Part~0 terms are unchanged; the
new term is zero for equivariant nominal models by construction.}
\label{tab:degradation2}
\end{table}

\begin{figure}[ht]
\centering
\resizebox{\textwidth}{!}{%
\begin{tikzpicture}[
  font=\footnotesize,
  box/.style={draw, rounded corners, align=center, minimum height=8mm, inner sep=3pt, text width=40mm},
  learn/.style={box, fill=orange!20, draw=orange!70!black, thick, text width=46mm},
  zero/.style={box, fill=blue!7},
  data/.style={box, fill=gray!12},
  arr/.style={-latex, thick},
  node distance=7mm and 6mm]
  \node[data] (tele) {telemetry $(x,u,y)$\\ phase registration $\to Z_k$};
  \node[learn, right=of tele] (model) {\textbf{nominal model $\hat f$}\\ analytic (Pinocchio observer, A1)\\ or learned: equivariant DeLaN\\ (mirror weight sharing, $\delta_f = 0$)};
  \node[data, right=of model] (res) {residual element $R_k = \mathfrak{r}(Z_k)$\\ $=\Pi^{+}R_k + \Pi^{-}R_k$ (+ base rows)};
  \node[zero, below=of model, xshift=-46mm] (rminus) {$R^{-}$: flip test on $\Pi^{-}R_k$\\ (Hemerik--Goeman, exact under $H_0$)};
  \node[zero, right=of rminus] (rplus) {joint rows: conformal on $\Pi^{+}R_k$\\ ($R^{+}$ magnitude, per-leg deviation)};
  \node[zero, right=of rplus] (base) {floating-base rows (6)\\ payload / base-level nuisances};
  \node[zero, below=of rplus, text width=70mm] (seq) {sequential layer: e-process / e-CUSUM (Ville)};
  \node[zero, below=of seq, text width=118mm] (iso) {isolation rules: pair--joint from the $R^{-}$ ranking $\cdot$ left/right from the joint rows $\cdot$ payload from the base rows};
  \node[data, left=4mm of rminus, text width=30mm] (mf) {model-free path $\hat f \equiv 0$:\\ $R_k = $ torque block of $Z_k$};
  \draw[arr] (tele) -- (model);
  \draw[arr] (model) -- (res);
  \draw[arr] (res.south) -- ++(0,-3.5mm) -| (rminus.north);
  \draw[arr] (res.south) -- ++(0,-3.5mm) -| (rplus.north);
  \draw[arr] (res.south) -- ++(0,-3.5mm) -| (base.north);
  \draw[arr] (rminus.south) -- ++(0,-3.5mm) -| (seq.west);
  \draw[arr] (rplus.south) -- (seq.north);
  \draw[arr] (base.south) -- ++(0,-3.5mm) -| (seq.east);
  \draw[arr] (seq) -- (iso);
  \draw[arr, dashed] (tele.south) -- ++(0,-3.5mm) -| (mf.north);
  \draw[arr, dashed] (mf.east) -- (rminus.west);
  \node[below=2mm of iso, align=center, text width=150mm] {\textcolor{orange!70!black}{orange: the only learned component (the nominal model)}\\ blue: zero trainable parameters (test, calibration, sequential layer, isolation rules)};
\end{tikzpicture}}
\caption{Three-channel architecture on the residual data element. The
nominal model is the only learned component (or the analytic observer); the
$R^{-}$ test reads $\Pi^{-}R_k$, the magnitude channel reads $\Pi^{+}R_k$ and
the joint rows, the base rows are the payload channel; the sequential and
isolation layers are fixed rules. The model-free $R^{-}$ of Sprints~1--2 is
the dashed path $\hat f \equiv 0$.}
\label{fig:architecture}
\end{figure}

\subsection{Empirical anchors (Sprint~6, e13; back-filled from the run)}
\label{sec:p2-anchors}

\begin{table}[p]
\centering
\footnotesize
\begin{tabular}{@{}>{\raggedright\arraybackslash}p{0.15\linewidth}>{\raggedright\arraybackslash}p{0.24\linewidth}p{0.55\linewidth}@{}}
\toprule
statement & experiment & status / reading \\
\midrule
Prop.~\ref{prop:residual-h0}, Cor.~\ref{cor:contamination} &
e13b: size of the residual flip test under $H_0$ vs $\delta_f^{(0.95)}$
(plain DeLaN ladder, equivariant ladder, analytic) &
e13b, $R=200$ nominal runs (go2\_urdf\_sym), flip test on the residual element, $\alpha=0.05$, band $[0.02,0.08]$: plain DeLaN residual (uncentred) size $1.00$ at every $\delta_f^{(0.95)} \in \{0.67, 1.47, 3.64, 14.8\}$\,N\,m for $K=60$ and $K=200$; equivariant DeLaN $0.030$--$0.050$ ($K=60$), $0.015$--$0.050$ ($K=200$); analytic residual $0.035/0.060$; raw $Z$ $0.015/0.035$ (URDF world conservative, D005). Naive centring: $1.00$ for every variant (Lemma~\ref{lem:centring}(iii)); $H_0'$ differenced test: $0.00$--$0.045$ for every variant incl.\ the plain $\delta_f^{(0.95)}=14.8$ model (Lemma~\ref{lem:centring}(iv)). Supplement (same seeds, $K\in\{5,10,20,60,200\}$): plain residual size $1.00$ from $K=10$ on for every $\delta_f$ of the ladder; naive centring $0.3/0.7/1.0/1.0$ at $K=10/20/60/200$ (growing with $K/K_{\mathrm{cal}}$ as in (iii)); differenced test in band at all $K$. \textbf{Confirmed} (contamination total already at $K=10$; equivariant/analytic in band or conservative). \\
Prop.~\ref{prop:power-gain} &
e13a: residual vs raw $R^{-}$ power on the low-SNR grid; variance
decomposition $\operatorname{Var}(\Pi^-y)$ vs $\operatorname{Var}(\Pi^-r)$;
Mahalanobis reference &
e13a, $R=50$, e-CUSUM at FAR $0.05$, alarm within 100 cycles: minimal detectable severity (det$\ge 0.9$) residual $R^-$ (analytic / equivariant DeLaN) vs raw $R^-$: bias $0.10$ vs $0.50$\,N\,m (both joints), gain $1-\kappa$ $0.02$ vs $0.05$ (HFE) and $0.02$ vs $0.02$ (KFE), friction $\times 1.5$ vs undetected on the grid; nominal $\operatorname{Var}(\Pi^-r)/\operatorname{Var}(\Pi^-\tau_{\mathrm{cmd}})$ median over rows $0.16$ (analytic), $0.47$ (DeLaN); fault SNR ratio $>1$ in all 22 cells for the analytic residual and in all but the three KFE-gain cells for the DeLaN residual (response term, Remark~\ref{rem:residual-weaker}(i)); Mahalanobis magnitude reference: bias $0.05$, gain $0.01$, friction $\times 1.5$. \textbf{Confirmed} for (a)+(b); the KFE-gain DeLaN cells are the conditional case (c). \\
Cor.~\ref{cor:two-channel} + Thm.~\ref{thm:n1-2} on residuals &
e13c: e04c groups (single / bilateral equal / unequal) with $R^{-}$ on
$\Pi^{-}r$ and $R^{+}$ on $\Pi^{+}r$; three-channel isolation confusion &
e13c, $R=50$, $\kappa$ groups of e04c on the equivariant residual: single --- $R^-$ on $\Pi^-r$ power $1.00$ (e-CUSUM median delay 15 cycles), $R^+$ on $\Pi^+r$ $1.00$ (3 cycles); bilateral equal --- $R^-$ e-process $0.00$, e-CUSUM $0.04$ (FAR level), $\Pi^-$-energy score $0.00$, $\Pi^+$ score $1.00$; unequal --- $R^-$ $1.00$; antisymmetric energy share of the residual signature $0.49 / 0.000 / 0.15$ (raw $Z$: $0.43/0.001/0.06$; a one-joint footprint splits exactly $1/2$--$1/2$). Three-channel isolation ($R=20$, 7 classes incl.\ lateral payload): raw+tracking $0.57$, analytic rows $0.86$, equivariant rows $1.00$. \textbf{Confirmed} (Theorem~\ref{thm:n1-2} verbatim on $R_k$). \\
Prop.~\ref{prop:equiv-delan} &
Block~Q: $\delta_f^{(0.95)}$ and validation RMSE, plain vs equivariant
ladders &
Block~Q (rp012): $\delta_f^{(0.95)}$ of the plain per-leg DeLaN $0.67 / 1.47 / 3.64 / 14.8$\,N\,m at $400$k$/50$k$/10$k$/2$k samples per leg ($\delta_f^{(0.5)}$ $0.21/0.40/1.28/3.44$); mirror-shared model $0$ to the last bit; validation RMSE equivariant vs plain $0.351/0.415$, $0.539/0.559$, $0.80/1.08$, $3.04/2.42$ (2k: 40 gradient steps, seed spread $\pm 0.6$--$0.8$); single template $0.325$; welded trunk $0.573/0.625$. \textbf{Confirmed} ($\delta_f=0$; no accuracy cost except the under-trained 2k cell). \\
\bottomrule
\end{tabular}
\caption{Empirical anchors of Part~2, filled from the e13 run
\texttt{e13-20260815-2334} (rp013) and the Block~Q reports (rp012);
wheeled-M1 anchors (e13d) are pending the M1 world.}
\label{tab:anchors}
\end{table}
